\subsection{Quellenumwandlung}
\Aufgabe{
Gegeben seien die zwei abgebildeten Netzwerke eines Generators und eines Verbrauchers.
Der Generator besteht aus zwei parallel geschalteten, realen Quellen mit den idealen
Spannungsquellen $U_{01}$ und $U_{02}$ sowie den Innenwiderständen $R_{i1}$ und $R_{i2}$.
Der Verbraucher besteht nur aus dem Lastwiderstand $R_V$. \newline
    
    \begin{center}
    \begin{tikzpicture}
 % --- Generator (linke Grafik) ---
    % äußerer Rahmen
    \draw[dashed] (0,-2.7) rectangle (5,2.7);

    % Klemmenleitungen
    \draw (1, 2) -- (4, 2);
    \draw (1,-2) -- (4,-2);

    % Klemmenpunkte und Bezeichnungen
    \node[ocirc,label=right:$A_C$] (AC) at (4, 2) {};
    \node[ocirc,label=right:$B_C$] (BC) at (4,-2) {};

    % Klemmenspannung U_G
    \draw[->] (4.2, 1.6) -- (4.2,-1.6) node[midway,right] {$U_G$};

    % erste reale Quelle: U01 + Ri1
    \draw[dashed] (0.5,-2.3) rectangle (1.5,2.3);
    \draw (1,2) to[V,v, name=U1] (1,0)
              to[R,l=$R_{i1}$, name=R1] (1,-2);

      % zweite reale Quelle: U02 + Ri2
    \draw[dashed] (2.5,-2.3) rectangle (3.5,2.3);
    \draw (3,2) to[V,l=$U_{02}$,i>_=$I_2$] (3,0)
              to[R,i,l=$R_{i2}$, name=R2] (3,-2);

    % --- Verbraucher (rechte Grafik) ---
    % äußerer Rahmen
    \draw[dashed] (5.6,-2.7) rectangle (8.8,2.7);

    % Klemmenpunkte
    \node[ocirc,label=left:$A_V$] (AV) at (6.3, 2) {};
    \node[ocirc,label=left:$B_V$] (BV) at (6.3,-2) {};

    % Verbraucherzweig
    \draw (AV) -- (8,2)
          to[R,l=$R_V$,i>_=$I_V$] (8,-2)
          -- (BV);

    \varrmore{U1}{$U_{\mathrm{01}}$};
    \iarrmore{R1}{$I_{\mathrm{1}}$};
  %      \varrmore{U1}{$U_{\mathrm{01}}$};
  %  \iarrmore{R1}{$I_{\mathrm{1}}$};
    
\end{tikzpicture}
    \end{center}
    \begin {enumerate}[label=\alph*)]
        \item Wie groß ist die Klemmenspannung $U_G$ des Generators?
             Wie groß sind die Ströme $I_1$ und $I_2$ im Generator?
             Welche Leistung $P_G$ wird im Generator umgesetzt?
             Begründen Sie Ihre Antwort mathematisch mit Hilfe von Maschen-
             und Knotengleichungen!
        \item Wie groß wäre die im Generator umgesetzte Leistung $P_G$
            in Abhängigkeit der gegebenen Größen $U_0$ und $R_i$, wenn
            aufgrund einer Fertigungstoleranz die Quellspannung $U_{01}$
            um $10\%$ größer als die Quellspannung $U_{02}$ bzw.\ $U_0$ wäre?
            Wie groß wäre in diesem Fall die Klemmenspannung $U_G$
            in Abhängigkeit von $U_0$?

        Für die weiteren Teilaufgaben soll der Generator mit dem Verbraucher $R_V$
        belastet werden. Es gelte weiter $U_0 = U_{01} = U_{02}$ und
        $R_i = R_{i1} = R_{i2}$.


        \item Wie groß ist jetzt die Klemmenspannung $U_G$ in Abhängigkeit
            von $U_0$, $R_i$ und $R_V$?
        \item Wie groß sind die im Generator umgesetzte Leistung $P_G$ und
            die im Verbraucher umgesetzte Leistung $P_V$ in Abhängigkeit
            von $U_0$, $R_i$ und $R_V$?
    \end{enumerate}
}

\Loesung{


\begin {enumerate}[label=\alph*)]
        \item Wie groß ist die Klemmenspannung $U_G$ des Generators?
             Wie groß sind die Ströme $I_1$ und $I_2$ im Generator?
             Welche Leistung $P_G$ wird im Generator umgesetzt?
             Begründen Sie Ihre Antwort mathematisch mit Hilfe von Maschen-
             und Knotengleichungen!

             \textbf{Maschen- und Knotengleichungen:}

\[
\sum_k U_k = 0:\quad
0 = U_{01} - R_{i1} I_1 - U_{02} + R_{i2} I_2
\]

\[
\sum_k I_k = 0:\quad
0 = I_1 + I_2
\;\Rightarrow\;
I_1 = -I_2
\]

\medskip
\textbf{Ströme $I_1$ und $I_2$:}

\[
\Rightarrow\; U_{01} - U_{02} = (R_{i1} + R_{i2}) I_1 = 0
\]

\[
\Rightarrow\; I_1 = -I_2 = 0
\]

\medskip
\textbf{Leistung $P_G$:}

\[
\Rightarrow\; P_G = R_{i1} I_1^2 + R_{i2} I_2^2
\quad \text{mit } R_i = R_{i1} = R_{i2} \text{ und } I_1 = -I_2 = 0
\]

\[
= 2 R_i I_1^2 = 0
\]

\medskip
\textbf{Klemmenspannung $U_G$:}

\[
\Rightarrow\; U_G = U_{01} - R_{i1} I_1 \quad \text{mit } I_1 = 0
\]

\[
= U_0 \quad \text{mit } U_0 = U_{01}
\]





        \item Wie groß wäre die im Generator umgesetzte Leistung $P_G$
            in Abhängigkeit der gegebenen Größen $U_0$ und $R_i$, wenn
            aufgrund einer Fertigungstoleranz die Quellspannung $U_{01}$
            um $10\%$ größer als die Quellspannung $U_{02}$ bzw.\ $U_0$ wäre?
            Wie groß wäre in diesem Fall die Klemmenspannung $U_G$
            in Abhängigkeit von $U_0$?

            \[
\Rightarrow\; U_{01} - U_{02} = (1{,}1 - 1)\, U_0 = 2 R_i I_1
\]

\[
\Rightarrow\; I_1 = -I_2 = \frac{0{,}1\, U_0}{2 R_i} = \frac{1}{20 R_i}\, U_0
\]

\medskip
\textbf{Leistung $P_G$:}

\[
\Rightarrow\; P_G = 2 R_i \left( \frac{1}{20 R_i} U_0 \right)^2
\]

\[
= \frac{1}{200 R_i}\, U_0^2
\]

\medskip
\textbf{Klemmenspannung $U_G$:}

\[
\Rightarrow\; U_G = U_{01} - R_{i1} I_1
\]

\[
= 1{,}1\, U_0 - R_i \frac{1}{20 R_i} U_0
\quad \text{mit } R_{i1} = R_i
\]

\[
= 1{,}05\, U_0
\]


        Für die weiteren Teilaufgaben soll der Generator mit dem Verbraucher $R_V$
        belastet werden. Es gelte weiter $U_0 = U_{01} = U_{02}$ und
        $R_i = R_{i1} = R_{i2}$.


        \item Wie groß ist jetzt die Klemmenspannung $U_G$ in Abhängigkeit
            von $U_0$, $R_i$ und $R_V$?

    \begin{center}
    \begin{tikzpicture}
    % --- Generator (linke Grafik) ---
    % äußerer Rahmen
    %\draw[dashed] (0,-2.7) rectangle (5.1,2.7);

    % Klemmenleitungen (Generator-Sammelschienen)


    \draw (6, 2)to[short] (7.5,2);
    \draw (6, -2)to[short] (7.5,-2);
    
    \draw (1, 2)to[short,-o] (6,2);
    \draw (1, -2)to[short,-o] (6,-2);

    % Klemmenpunkte und Bezeichnungen
    %\node[ocirc,label=right:$A_\mathrm{G}$] (AC) at (4, 2) {};
    %\node[ocirc,label=right:$B_\mathrm{G}$] (BC) at (4,-2) {};

    % Klemmenspannung U_G
    \draw[->] (6, 1.6) -- (6,-1.6) node[midway,right] {$U_G$};

    % --- 1. Quelle: Norton (reale Stromquelle) I01 || Ri1 ---
    \draw[dashed] (0.2,-1.6) rectangle (2.6,1.6);

    % Abzweigpunkte von den Sammelschienen zur Quelle
    \draw (1,2) -- (0.8,2);
    \draw (1,-2) -- (0.8,-2);
    \draw (1,2) -- (1.2,2);
    \draw (1,-2) -- (1.2,-2);

    % parallele Zweige: Stromquelle und Innenwiderstand
    \draw (0.8,2)  to[I,i>, name=I1] (0.8,-2);
    \draw (1.8,2)  to[R,l=$R_\mathrm{i1}$]   (1.8,-2);

    % --- 2. Quelle: Norton (reale Stromquelle) I02 || Ri2 ---
    \draw[dashed] (2.85,-1.6) rectangle (5.3,1.6);

    % Abzweigpunkte von den Sammelschienen zur Quelle
    \draw (3,2) -- (2.8,2);
    \draw (3,-2) -- (2.8,-2);
    \draw (3,2) -- (3.2,2);
    \draw (3,-2) -- (3.2,-2);

    % parallele Zweige: Stromquelle und Innenwiderstand
    \draw (3.4,2)  to[I,i>, name=I2] (3.4,-2);
    \draw (4.4,2)  to[R,l=$R_\mathrm{i2}$]   (4.4,-2);

    % --- Verbraucher (rechte Grafik) ---
    % äußerer Rahmen
    %\draw[dashed] (5.6,-2.7) rectangle (8.8,2.7);



    % Verbindung Generator -> Verbraucher
    %\draw (AC) -- (AV);
    %\draw (BC) -- (BV);

    \node (AV) at (7.5,2) {}; 
    \node (BV) at (7.5,-2) {};

    % Verbraucherzweig
    \draw (AV) -- (7.5,2)
          to[R,l=$R_V$] (7.5,-2)
          -- (BV);

  %  \iarrmore{RV}{$I_\mathrm{V}$};
     \iarrmore{I1}{$I_\mathrm{01}$};
     \iarrmore{I2}{$I_\mathrm{02}$};
\end{tikzpicture}

    \end{center}

    Anschließend die beiden realen Stromquellen zusammenfassen und zu einer realen Spannungsquelle umrechnen: 

       \begin{center}
    \begin{tikzpicture}
    % Klemmenleitungen / Klemmen
    \draw (4, 2)to[short] (5.5,2);
    \draw (4,-2)to[short] (5.5,-2);

    \draw (1, 2)to[short,-o] (4,2);
    \draw (1,-2)to[short,-o] (4,-2);

    % Klemmenspannung U_G
    \draw[->] (4, 1.6) -- (4,-1.6) node[midway,right] {$U_\mathrm{G}$};

    % --- Zusammenfassung: reale Spannungsquelle U0,ges in Serie mit Ri,ges ---
    \draw[dashed] (0.2,-1.8) rectangle (3.0,1.7);

    % Serienzweig: Spannungsquelle + Innenwiderstand
    \draw (1,2)
          to[V,v, name=Uges] (1,0)
          to[R,l=$R_{i,\mathrm{ges}}$, name=Rges] (1,-2);

    % Beschriftung U0,ges an der Quelle
    \varrmore{Uges}{$U_{0,\mathrm{ges}}$};

    % --- Verbraucherzweig ---
    \draw (5.5,2)
          to[R,l=$R_\mathrm{V}$] (5.5,-2);
        
\end{tikzpicture}
    \end{center}

    \[
I_{01} = \frac{U_{01}}{R_{i1}} = \frac{U_0}{R_i}
\]

\[
I_{02} = \frac{U_{02}}{R_{i2}} = \frac{U_0}{R_i}
\]

\medskip
\textbf{Zusammenfassen von Quellen und Widerständen:}

\[
I_{0,\mathrm{ges}} = I_{01} + I_{02}
= 2\,\frac{U_0}{R_i}
\]

\[
R_{i,\mathrm{ges}}
= \frac{R_{i1} R_{i2}}{R_{i1} + R_{i2}}
= \frac{R_i^2}{2 R_i}
= \frac{1}{2} R_i
\]

\textbf{Stromquelle in äquivalente Spannungsquelle umwandeln:}

\[
U_{0,\mathrm{ges}}
= R_{i,\mathrm{ges}}\, I_{0,\mathrm{ges}}
= \frac{1}{2} R_i \, 2 \frac{U_0}{R_i}
= U_0
\]

\medskip
\textbf{Generatorspannung $U_G$:}

\[
U_G
= \frac{R_V}{R_V + R_{i,\mathrm{ges}}}\, U_{0,\mathrm{ges}}
\]

\[
= \frac{R_V}{R_V + \tfrac{1}{2} R_i}\, U_0
\]

\[
= \frac{2 R_V}{2 R_V + R_i}\, U_0
\]


        \item Wie groß sind die im Generator umgesetzte Leistung $P_G$ und
            die im Verbraucher umgesetzte Leistung $P_V$ in Abhängigkeit
            von $U_0$, $R_i$ und $R_V$?


            \[
P_G
= \frac{1}{R_{i,\mathrm{ges}}}\,\bigl(U_{0,\mathrm{ges}} - U_G\bigr)^2
\]

\[
= \frac{2}{R_i}
\left(
U_0 - \frac{2 R_V}{2 R_V + R_i}\, U_0
\right)^2
\]

\[
= \frac{2}{R_i}
\left(
1 - \frac{2 R_V}{2 R_V + R_i}
\right)^2
U_0^2
\]

\medskip

\[
P_V
= \frac{1}{R_V}\, U_G^2
\]

\[
= \frac{1}{R_V}
\left(
\frac{2 R_V}{2 R_V + R_i}
\right)^2
U_0^2
\]
    \end{enumerate}

}